(1) is the conjunction of Proposition \claimref{families.three-identical-families-mixing-shared-slots} (the faces are $S_{3}$-identical), the definition of the bracket as face-valued (one bracket per two-element subset of slots, Fact \claimref{families.imported-face-valued-bracket-quantum-numbers}), and Proposition \claimref{families.alignment-computed-thirteenth-refutation-democracy-one}(2)-(3) (the vacuum value has a direction in face space and the direction orders the faces by overlap); it is structural because each ingredient is. (2) The texture argument is the flavour literature's \cite{GST1968, Fritzsch1978}: a nearest-neighbour Yukawa matrix in the down sector with entries $(0, A; A, B)$ on the light two-by-two block has eigenvalues of ratio $m_{d}/m_{s} \approx A^{2}/B^{2}$ and mixing angle $\tan\theta_{C} \approx A/B$, whence $\theta_{C} \approx \sqrt{m_{d}/m_{s}}$; the assumptions (hierarchy, a vanishing $(1,1)$ entry, the up sector's contribution neglected) are the inherited ones. The arithmetic: with $m_{d} = 4.70$ MeV and $m_{s} = 93.5$ MeV \cite{PDG2024} ($\overline{\mathrm{MS}}$ at $2$ GeV), $\sqrt{m_{d}/m_{s}} = 0.2242$ against $0.225$, agreement to $0.4$ percent on central values, well inside the few-percent band the quark-mass uncertainties allow; the script verifies $0.215^{2} < m_{d}/m_{s} < 0.235^{2}$ as exact inequalities.
